\documentclass{ctexart}
\usepackage{template}
\usepackage{esint,extarrows,tikz}
\begin{document}\pagestyle{empty}
\begin{center}\Large
北京大学数学科学学院2023-24高等数学A1期末考试
\end{center}
\begin{problem}[1.(20\songti{分})]
求下列函数的极限.
        \begin{enumerate}[label=\tbf{(\arabic*)}]
            \item \textbf{(10\songti{分})}\ 求极限$\displaystyle\lim_{x\to0}\dfrac1x\left(\left(1+x\right)^{\frac1x}-\e\right)$.
            \item \textbf{(10\songti{分})}\ 设函数$f(x)$在$x=0$处$n+1$阶可导,且满足\[f(0)=f'(0)=\cdots=f^{(n-1)}(0)=0\ \ \ \ \ f^{(n)}(0)=a\]求极限$\displaystyle\lim_{x\to0}\dfrac{f\left(\e^x-1\right)-f(x)}{x^{n+1}}$.
        \end{enumerate}
\end{problem}
\begin{solution}
\begin{enumerate}[label=\tbf{(\arabic*)},leftmargin=*]
\item 由$\ln(1+x)/x=1-x/2+O(x^2)$,有
\[(1+x)^{1/x}=\e\exp(-x/2+O(x^2))=\e-\dfrac\e2x+O(x^2)\]
故所求极限为$-\e/2$.
\item 设$f^{(n+1)}(0)=b$,由Peano型Taylor公式,
\[f(t)=\dfrac{a}{n!}t^n+\dfrac{b}{(n+1)!}t^{n+1}+o(t^{n+1})\]
又$\e^x-1=x+x^2/2+O(x^3)$,故
\[(\e^x-1)^n=x^n+\dfrac n2x^{n+1}+o(x^{n+1}),\qquad
(\e^x-1)^{n+1}=x^{n+1}+o(x^{n+1})\]
于是$f(\e^x-1)-f(x)=\dfrac{a}{2(n-1)!}x^{n+1}+o(x^{n+1})$,极限为$\dfrac{a}{2(n-1)!}$.
\end{enumerate}
\end{solution}
\begin{problem}[2.(20\songti{分})]
回答下列问题.
        \begin{enumerate}[label=\tbf{(\arabic*)}]
            \item \textbf{(10\songti{分})}\ 设函数$F(u,v)$有连续的二阶偏导数,$z=z(x,y)$是由方程$F(x-z,y-z)=0$确定的隐函数.计算并化简
                \[\dfrac{\p^2z}{\p x^2}+\dfrac{\p^2z}{\p x\p y}+\dfrac{\p^2z}{\p y\p x}+\dfrac{\p^2z}{\p y^2}\]
            \item \textbf{(10\songti{分})}\ 给定方程组
                \[\left\{\begin{array}{l}
                    xy+yz^2+4=0\\
                    x^2y+yz-z^2+5=0
                \end{array}\right.\]
                试讨论上述方程在$P_0\left(1,-2,1\right)$处能确定的隐函数,并计算其在$P_0$处的导数.
        \end{enumerate}
\end{problem}
\begin{solution}
\begin{enumerate}[label=\tbf{(\arabic*)},leftmargin=*]
\item 在通常的隐函数非退化条件$F_u+F_v\neq0$下,
\[z_x=\dfrac{F_u}{F_u+F_v},\qquad z_y=\dfrac{F_v}{F_u+F_v}\]
所以$z_x+z_y=1$,分别对$x,y$求导并相加,得到所求四项之和为$0$.
现有题干未明写非退化条件;若允许$F\equiv0$,则任意$z(x,y)$都满足方程,结论不再确定.因此该条件为通常隐函数解释下的必要补充.
\item 在$P_0=(1,-2,1)$处,
\[\nabla F=(-2,2,-4),\qquad\nabla G=(-4,2,-4)\]
相应的雅可比子式为
\[\dfrac{D(F,G)}{D(x,y)}=4,\qquad
\dfrac{D(F,G)}{D(x,z)}=-8,\qquad
\dfrac{D(F,G)}{D(y,z)}=0\]
故可将$x,y$表示为$z$的光滑函数,求导得到
\[-2x'+2y'-4=0,\qquad -4x'+2y'-4=0\]
即$\dfrac{\di x}{\di z}=0$, $\dfrac{\di y}{\di z}=2$.
也可将$x,z$表示为$y$的函数,对应导数为$\dfrac{\di x}{\di y}=0$, $\dfrac{\di z}{\di y}=1/2$.
不能在该点将$y,z$同时表示为$x$的可微函数,否则两方程求导后相减将得到$2=0$.
\end{enumerate}
\end{solution}
\begin{problem}[3.(20\songti{分})]
求函数$f(x,y)=\left(y-x^2\right)\left(y-x^3\right)$的极值.
\end{problem}
\begin{solution}
令$f_x=f_y=0$,由
\[f_x=5x^4-3x^2y-2xy,\qquad f_y=2y-x^2-x^3\]
得到驻点$(0,0),(1,1),(2/3,10/27)$.
记$A=f_{xx}=20x^3-6xy-2y$, $B=f_{xy}=-3x^2-2x$, $C=f_{yy}=2$.
在$(0,0)$附近,$f(x,0)=x^5$变号,故不是极值点.
在$(1,1)$处,$AC-B^2=-1<0$,故为鞍点.
在$(2/3,10/27)$处,$A=100/27>0$, $B=-8/3$, $AC-B^2=8/27>0$,故为严格极小值点,极小值为
\[f(2/3,10/27)=-\dfrac4{729}\]
函数没有其他极值点,也没有极大值.
\end{solution}
\begin{problem}[4.(20\songti{分})]
回答下列问题.
        \begin{enumerate}[label=\tbf{(\arabic*)}]
            \item \textbf{(10\songti{分})}\ 设函数$f(x,y)$在点$(0,0)$的某邻域内有定义且在$(0,0)$处连续.若极限$\displaystyle\lim_{(x,y)\to(0,0)}\dfrac{f(x,y)}{x^2+y^2}$存在,试证明$f(x,y)$在$(0,0)$处可微.
            \item \textbf{(10\songti{分})}\ 欧氏空间$\R^3$中平面$T:x+y+z=1$截圆柱面$S:x^2+y^2=1$得一椭圆周$R$.求$R$上到原点最近和最远的点.
        \end{enumerate}
\end{problem}
\begin{solution}
\begin{enumerate}[label=\tbf{(\arabic*)},leftmargin=*]
\item 记该有限极限为$L$,则$f(x,y)=(x^2+y^2)(L+o(1))$.由连续性,$f(0,0)=0$.
令$\rho=\sqrt{x^2+y^2}$,则$f(x,y)/\rho=\rho(L+o(1))\to0$.
因此$f$在原点可微,且全微分为零.
\item 在椭圆上有$x^2+y^2=1$,故到原点的距离平方为$1+z^2$,且$z=1-x-y$.
由于$-\sqrt2\leqslant x+y\leqslant\sqrt2$,可取到$z=0$,对应最近点为$(1,0,0),(0,1,0)$,距离为$1$.
最大的$|z|$为$1+\sqrt2$,对应最远点
\[\left(-\dfrac1{\sqrt2},-\dfrac1{\sqrt2},1+\sqrt2\right)\]
其距离为$\sqrt{4+2\sqrt2}$.
\end{enumerate}
\end{solution}
\begin{problem}[5.(20\songti{分})]
回答下列问题.
        \begin{enumerate}[label=\tbf{(\arabic*)}]
            \item \textbf{(10\songti{分})}\ 设$f(x)$是一个定义在$\R$上的周期为$T\neq0$的无穷阶光滑函数.试证明:对于任意$k\in\N^*$,总存在$\xi\in\R$使得$f^{(k)}(\xi)=0$.
            \item \textbf{(10\songti{分})}\ 设函数$f(u,v)$有连续的偏导数$f_u(u,v)$和$f_v(u,v)$且满足$f(x,1-x)=1$.试证明:在单位圆周$S:u^2+v^2=1$上至少存在两个不同的点$(u_1,v_1)$和$(u_2,v_2)$使得$v_if_u(u_i,v_i)=u_if_v(u_i,v_i)$,其中$i=1,2$.
        \end{enumerate}
\end{problem}
\begin{solution}
\begin{enumerate}[label=\tbf{(\arabic*)},leftmargin=*]
\item 可将周期取为$|T|>0$.各阶导数也以$|T|$为周期.
对任意$k\geqslant1$,函数$f^{(k-1)}$在任一长度为$|T|$的闭区间两端取值相等,由Rolle定理,区间内部存在$\xi$使$f^{(k)}(\xi)=0$.
\item 令$g(\theta)=f(\cos\theta,\sin\theta)$,则
\[g'(\theta)=-\sin\theta f_u(\cos\theta,\sin\theta)+\cos\theta f_v(\cos\theta,\sin\theta)\]
题设给出$g(0)=g(\pi/2)=g(2\pi)=1$.在$[0,\pi/2]$和$[\pi/2,2\pi]$分别应用Rolle定理,得到位于两个互不相交开区间的$\theta_1,\theta_2$,使$g'(\theta_i)=0$.
对应的圆周点互异,且满足$v_if_u(u_i,v_i)=u_if_v(u_i,v_i)$.
\end{enumerate}
\end{solution}
\end{document}
